% EDIT HERE: title, author, affiliation, accent colour, then the sample content.
% All examples are self-contained: no images or bibliography tools are needed.
\documentclass[11pt,a4paper]{article}
\usepackage[T1]{fontenc}
\usepackage{tgpagella}
\usepackage[margin=25mm]{geometry}
\usepackage{microtype,amsmath,amssymb,graphicx,booktabs,tabularx}
\usepackage{xcolor,titlesec,fancyhdr,enumitem}
\definecolor{accent}{HTML}{245C61}
\definecolor{muted}{HTML}{58666B}
\usepackage[colorlinks=true,allcolors=accent]{hyperref}
\newcommand{\articletitle}{Article title}
\newcommand{\authorname}{Your name}
\newcommand{\affiliation}{Department / Institution}
\hypersetup{pdftitle={\articletitle},pdfauthor={\authorname}}
\titleformat{\section}{\large\bfseries\color{accent}}{\thesection}{0.7em}{}
\titleformat{\subsection}{\normalsize\bfseries}{\thesubsection}{0.7em}{}
\titlespacing*{\section}{0pt}{18pt}{7pt}
\setlength{\parindent}{0pt}
\setlength{\parskip}{6pt}
\setlist[itemize]{leftmargin=*,itemsep=2pt,topsep=3pt}
\pagestyle{fancy}\fancyhf{}
\fancyhead[L]{\small\color{muted}\authorname}
\fancyhead[R]{\small\color{muted}Research article}
\fancyfoot[C]{\small\thepage}
\renewcommand{\headrulewidth}{0pt}
\setlength{\headheight}{15pt}

\begin{document}
\thispagestyle{plain}
{\small\sffamily\color{accent}RESEARCH ARTICLE}\par
\vspace{8pt}
{\LARGE\bfseries\articletitle\par}
\vspace{7pt}
{\large\authorname}\par
{\small\color{muted}\affiliation\quad /\quad \today}\par
\vspace{6pt}{\color{accent}\hrule height 1.2pt}\vspace{10pt}
\textbf{Abstract}\quad
State the problem, the approach and the main finding in a single paragraph.
This example describes an evaluation of a research workflow: a baseline process
is compared with a revised process using completion time and error rate.
Replace this summary with the contribution of your own study and its implications.

{\small\textbf{Keywords:} research design; evaluation; reproducibility}
\section{Introduction}
Open with the problem and explain who is affected by it. Establish what prior
work already tells us, what remains uncertain and why resolving that gap matters.
Use a reference where a claim needs support, for example \cite{example}.

The objective of this study is to evaluate whether a revised workflow improves
performance without sacrificing accuracy. The example numbers below illustrate
the layout; they are not research findings.
\subsection{Research question and contribution}
\begin{itemize}
\item \textbf{Question:} How does the revised workflow compare with the baseline?
\item \textbf{Contribution:} A reproducible comparison using two explicit metrics.
\item \textbf{Scope:} One task, two conditions and a clearly defined sample.
\end{itemize}
\section{Methods}
\subsection{Study design}
Describe the participants or materials, the procedure and the comparison groups.
Specify inclusion criteria and explain how you selected the sample. Report enough
detail for another researcher to repeat the work.
\subsection{Measures and analysis}
Define your primary outcome before presenting results. For example, the relative
improvement in completion time is
\begin{equation}
\Delta_t = \frac{t_{\mathrm{baseline}}-t_{\mathrm{revised}}}{t_{\mathrm{baseline}}}\times100\%.
\label{eq:improvement}
\end{equation}
Report uncertainty alongside averages and state how missing observations were
handled. Separate planned analyses from exploratory ones.
\section{Results}
Table~\ref{tab:results} shows an example comparison. Replace the illustrative values,
units and measures with your own results. Equation~\ref{eq:improvement} gives a
relative reduction of 25\% for the example completion times.
\begin{table}[htbp]
\centering
\caption{Illustrative outcomes; replace with measured data.}
\label{tab:results}
\begin{tabularx}{\linewidth}{@{}Xrr@{}}
\toprule
\textbf{Measure} & \textbf{Baseline} & \textbf{Revised}\\
\midrule
Completion time (minutes) & 24.0 & 18.0\\
Error rate (\%) & 8.0 & 5.0\\
Completed tasks & 40 & 40\\
\bottomrule
\end{tabularx}
\end{table}
\section{Discussion}
Interpret the finding in relation to the research question rather than repeating
the table. Compare it with prior work and explain which mechanisms could account
for the observed difference.
\subsection{Limitations}
Discuss sample size, measurement assumptions and threats to validity. Explain
which settings the result may generalize to and where further evidence is needed.
\section{Conclusion}
Give a direct answer to the research question, state the main contribution and
identify one concrete next step. Avoid introducing new evidence here.
\section*{Data and acknowledgements}
Describe where materials and analysis code can be found, any access restrictions,
funding and contributions that do not meet the authorship criteria.
\begin{thebibliography}{9}
\bibitem{example} A. Author. \textit{Title of a relevant work}. Publisher, 2025.
% Replace this illustrative entry with a real source. Add entries with \bibitem.
\end{thebibliography}
\end{document}
